\section{Liquid metal in a magnetic field: MHD in FluidX3D}\label{sec:mhd}

In the section on liquid lead (Section~\ref{sec:filament}) the pressure drop in the magnetic field was only estimated
with a rule of thumb: $\mathrm{d}p/\mathrm{d}x = \sigma\,v\,B^2 k_p$, with $k_p \approx 1/\mathrm{Ha}$ for insulating
walls. The user asked whether the flow in the magnetic field could also be computed dynamically, as a kind of CFD
simulation, and whether this could be coupled with FluidX3D. FluidX3D has no magnetic fields of its own. This section
describes how the required force was added in two stages and checked against cases with known solutions.

\paragraph{What happens physically?} When an electrically conducting metal moves across a magnetic field, a voltage
arises in it, as in a dynamo. It drives currents $\mathbf{j}$, and these currents in turn feel a force in the field,
the \emph{Lorentz force} $\mathbf{j}\times\mathbf{B}$. It brakes the motion across the field. At a small magnetic
Reynolds number, as for lead in a coil, the currents hardly change the external field (\emph{inductionless}
approximation). Then Ohm's law for moving conductors and charge conservation apply:
\[
\mathbf{j} = \sigma\left(-\nabla\varphi + \mathbf{u}\times\mathbf{B}\right),\qquad \nabla\cdot\mathbf{j} = 0
\quad\Rightarrow\quad \nabla^2\varphi = \nabla\cdot(\mathbf{u}\times\mathbf{B}),\qquad
\mathbf{f} = \frac{\mathbf{j}\times\mathbf{B}}{\rho}.
\]
$\sigma$ is the electrical conductivity, $\varphi$ the electric potential. How strongly the field acts compared with
viscosity is measured by the \emph{Hartmann number} $\mathrm{Ha} = B\,a\sqrt{\sigma/(\rho\nu)}$ ($a$ half the duct
width along the field). Thin boundary layers of thickness $a/\mathrm{Ha}$ form at the walls perpendicular to the field
(\emph{Hartmann walls}); thicker layers of thickness $a/\sqrt{\mathrm{Ha}}$ form at the walls parallel to the field
(\emph{side walls}) \cite{huntshercliff1971}. Whether the currents close through the walls or only through these
layers decides the pressure drop.

\paragraph{Terms used in this section.}
\begin{description}\itemsep2pt
\item[Hartmann walls] the two duct walls that the magnetic field meets at right angles. If the field points upwards,
  they are the floor and the ceiling of the duct. Named after the Danish physicist Julius Hartmann, who studied this
  flow in the 1930s.
\item[side walls] the other two walls, along which the field runs (left and right).
\item[insulating wall] a wall through which no electric current can flow, for example one with a ceramic layer. The
  currents generated in the metal must then close inside the metal, in the thin layers at the wall.
\item[conducting wall] a wall that conducts current well, for example copper or thick steel. The currents can then
  close through the wall as through a wire; the circuit has little resistance and the braking force becomes large.
  \emph{Perfectly conducting} means without any resistance, a limiting case for the calculation.
\item[conducting Hartmann walls] floor and ceiling conduct, the side walls insulate. This is Hunt's case (1965), for
  which an exact solution exists.
\item[Hartmann number Ha] a measure of how strongly the magnetic field acts compared with the viscosity of the metal.
  Ha = 0 means no field; at Ha = 100 the field almost completely dominates the flow.
\item[Hartmann layer] the thin layer at the Hartmann walls in which the velocity rises from zero to the value in the
  middle. Its thickness is half the duct width divided by Ha, i.e.\ one hundredth at Ha = 100.
\item[side layer] the corresponding layer at the side walls. It is thicker: half the duct width divided by
  $\sqrt{\mathrm{Ha}}$.
\item[potential $\varphi$] the electric voltage at each point in the metal, measured against a fixed point. Its
  gradient steers the currents so that they close.
\item[Jacobi step] a simple method for computing $\varphi$: each cell repeatedly takes the average of its neighbours,
  corrected by the source, until nothing changes any more.
\item[pressure-drop coefficient $k_p$] the fraction of the largest possible magnetic braking force that actually acts
  (1 = full braking, small = the currents find only a difficult return path).
\end{description}

\subsection*{Stage 1: Lorentz force without potential}
Between two infinitely wide plates $\nabla\varphi$ vanishes, leaving $\mathbf{f} = -K\,\mathbf{u}_\perp$ with
$K = \sigma B^2/\rho$. This force was built into the main kernel of the project's own FluidX3D copy, semi-implicitly
with the Guo forcing: $\mathbf{u}_\perp = (\mathbf{u}^*_\perp + \mathbf{F}_\perp/2\rho)/(1+K/2)$. This keeps the
calculation stable even for large $K$. The changes are published as diffs in the folder \texttt{fluidx3d/}
(FluidX3D licence \cite{fluidx3d}). The check is the Hartmann flow, which has an exact solution:
$u(y) = f_x/K\,(1-\cosh(\mathrm{Ha}\,y/a)/\cosh\mathrm{Ha})$.

\bild{mhd_hartmann_en.png}{1.0}{Stage 1: Hartmann flow between two plates. Left: FluidX3D (dots) against the exact
solution (lines) for Ha = 0 to 30. Right: the error at Ha = 30 versus grid resolution.}
\lesehilfe{On the left each colour is one field strength. Horizontally is the position between the two plates ($-1$
and $+1$ are the walls), vertically the velocity relative to its maximum. Without a field (grey) the familiar
parabola appears. The stronger the field, the flatter the profile in the middle and the steeper its drop at the
walls: the field brakes the fast middle, and only thin layers at the wall keep a gradient. The dots (FluidX3D) lie on
the lines (exact solution). On the right you see the error shrink as the grid is refined. The labels give the number
of cells inside the thin wall layer. Even with a single cell the error stays below 1\,\%; on finer grids it falls
roughly like the dashed line (second order).}

Without a field the deviation is 0.03\,\%, at Ha = 5 it is 0.04\,\%, and at Ha = 30 with 32 cells (just over one cell
in the Hartmann layer) 0.5\,\%. On finer grids the error falls roughly with second order (3.0 / 1.5 / 0.48 / 0.13\,\%
for $a$ = 8 / 16 / 32 / 64 cells).

\paragraph{Correction.} The values obtained first were much worse (4.3\,\% at Ha = 30 with 32 cells, 68\,\% with 8
cells) and already appeared in an earlier version. They were an error in the evaluation, not in the calculation.
Without the switch \texttt{UPDATE\_FIELDS}, FluidX3D recomputes the output velocity in a separate kernel, and that
kernel lacked half of the Lorentz force. The larger $K$, the larger the error. This only came to light when the
Hartmann test was repeated with the stage~2 code, in which \texttt{UPDATE\_FIELDS} is on. The rule ``at least two
cells in the Hartmann layer'' was therefore too strict; about one cell is enough for 1\,\% accuracy.

\subsection*{Stage 2: ducts with side walls}
In a real duct the currents flow in loops. In the middle they run across the flow; to return they must go through
the walls or through the thin layers at the walls. This return path is described by the potential $\varphi$.
Without $\varphi$, stage~1 would overestimate the braking force in a duct with insulating walls several times over.

\paragraph{Implementation.} A new switch \texttt{MHD\_POT} allocates two buffers for $\varphi$. Before every time
step the GPU performs some Jacobi steps for $\nabla\cdot\mathbf{j} = 0$, here 20 double steps. The equation is written
as a balance of currents over the six faces of each cell, so that charge is also conserved discretely. On an
insulating wall the current through the wall face is zero. A perfectly conducting wall holds $\varphi = 0$ on its
surface, which lies half a cell in front of the wall node. In the main kernel the force
$K\rho\,(\mathbf{E}\times\hat{\mathbf{B}})$ with $\mathbf{E} = -\nabla\varphi$ is added from the same face values.

\paragraph{Validation.} Comparing the code only with itself would be worthless. So a separate 2D solver for fully
developed duct flow was written that takes a different route. It does not use $\varphi$ but the induced magnetic
field $b$ (Shercliff \cite{shercliff1953}, Hunt \cite{hunt1965}):
\[
\nabla^2 u + \mathrm{Ha}\,\frac{\partial b}{\partial y} = -1,\qquad \nabla^2 b + \mathrm{Ha}\,\frac{\partial u}{\partial y} = 0 .
\]
It reproduces three known limits: the square duct without field ($\bar u = 0.1406$, i.e.\ $f\,\mathit{Re} = 14.23$),
the exact Hartmann solution for a wide duct with conducting walls, and the core velocity $1/\mathrm{Ha}$ for a wide duct
with insulating walls. Along the way a bug in its own boundary treatment was found and fixed.

Two classic cases in the square duct were checked. In \emph{Shercliff}'s case all walls are insulating. In
\emph{Hunt}'s case the Hartmann walls are perfectly conducting and the side walls insulating; the current returns
through the wall, and fast jets form at the side walls.

\paragraph{A bug that only the validation revealed.} In the first run the error with insulating walls grew strongly
with Ha: 12\,\% at Ha = 15, 36\,\% at Ha = 30. With conducting walls everything was fine. The reason: in the core of
an insulated duct two large forces almost cancel, the braking force $-K\mathbf{u}_\perp$ and the force from the
potential. Only a fraction of about $1/\mathrm{Ha}$ remains, so any small inconsistency between the two parts is
amplified by Ha. In the kernel the potential force was added only after the semi-implicit part, so that part missed
some of the force. After swapping the order, and using the same stored velocity for $\mathbf{u}\times\mathbf{B}$ as in
the Jacobi step, the error at Ha = 15 dropped from 12\,\% to 0.45\,\%.

\begin{table}[H]\centering\small
\begin{tabular}{lrrrr}\toprule
Case & Ha & $a$ [cells] & profile error & flow-rate error\\\midrule
no field & 0 & 32 & 0.01\,\% & $-0.03$\,\%\\
Shercliff (insulating) & 5 / 10 / 15 & 32 & 0.09 / 0.23 / 0.45\,\% & $-0.15$ / $-0.30$ / $-0.37$\,\%\\
Shercliff & 30 & 64 & 0.49\,\% & $-0.19$\,\%\\
Hunt (conducting Hartmann walls) & 5 / 10 / 15 & 32 & 0.08 / 0.42 / 0.86\,\% & $-0.09$ / $-0.30$ / $-0.68$\,\%\\
Hunt & 30 & 64 & 0.79\,\% & $-0.78$\,\%\\
grid study Shercliff & 10 & 16 / 32 / 64 & 0.75 / 0.23 / 0.09\,\% & \\\bottomrule
\end{tabular}\caption{Stage 2: FluidX3D against the independent 2D solver (RMS error of the profile relative to the
maximum; deviation of the mean flow rate).}\end{table}

\bild{mhd_kanal_en.png}{1.0}{Stage 2: square duct, B pointing straight up. Top left and centre: insulating walls
(Shercliff); top right: conducting Hartmann walls (Hunt). Bottom left and centre: the whole cross-section for Hunt,
Ha = 30, FluidX3D and reference. Bottom right: the grid study.}
\lesehilfe{You are looking into a square duct; the metal flows towards you and the magnetic field points up. Top
left: the velocity profile from bottom to top, i.e.\ along the field; top centre: from left to right, across the
field. Dots are FluidX3D, lines the independent reference solver; the percentages are the deviations. With
insulating walls a stronger field makes the profile flat and box-like along the field (thin layers at top and
bottom). Across the field it stays rounder because the side layers are thicker. Top right: the top and bottom walls
conduct. Now the field strongly brakes the middle, and the metal escapes into two fast jets along the side walls. The
two colour maps at the bottom show exactly this over the whole cross-section (yellow fast, dark slow); left FluidX3D,
centre the reference. They are hard to tell apart. Bottom right: the error falls as the grid is refined.}

All twelve cases agree with the reference solver to better than 1\,\%. The grid study shows an order of 1.4 to 1.7,
somewhat worse than in stage~1. The likely reason is the treatment of the wall faces; this has not been investigated
yet.

\paragraph{Checking the earlier rule of thumb.} The reference solver gives the pressure-drop coefficient
$k_p = \mathrm{d}p/\mathrm{d}x\,/\,(\sigma\bar u B^2) = 1/(\mathrm{Ha}^2\,\bar U)$ in the square duct:
\begin{table}[H]\centering\small
\begin{tabular}{rrrr}\toprule
Ha & $1/\mathrm{Ha}$ & $k_p$ insulating & $k_p$ conducting Hartmann walls\\\midrule
5 & 0.200 & 0.407 & 0.925\\
10 & 0.100 & 0.153 & 0.634\\
20 & 0.050 & 0.065 & 0.475\\
50 & 0.020 & 0.023 & 0.335\\
100 & 0.010 & 0.011 & 0.256\\\bottomrule
\end{tabular}\caption{Pressure-drop coefficient from the 2D solver (grid with at least three cells per Hartmann layer).}\end{table}
With insulating walls $k_p$ does approach $1/\mathrm{Ha}$; at Ha = 100 it is still 10\,\% above. The rule of thumb was
therefore right for large Ha and too low for small Ha. If the Hartmann walls conduct, $k_p$ falls only slowly, and at
Ha = 100 the pressure drop is about 25 times larger than with insulating walls.

\paragraph{Limits and next steps.}
\begin{itemize}\itemsep1pt
\item \textbf{Real Hartmann numbers are much larger.} Lead--lithium in the \SI{10}{mm} channels at \SI{6}{T} gives
  roughly Ha $\approx 600$, in blanket ducts of a few centimetres several thousand (material data not yet checked).
  With one cell per Hartmann layer, Ha = 600 alone would need about 600 cells per half duct width. Memory suffices,
  but the simple Jacobi solver converges too slowly on such fine grids. A multigrid method or a grid refined towards
  the wall would be needed; until then computed values hold up to Ha $\approx 30$, and the validated scaling laws
  beyond.
\item \textbf{Finitely conducting walls} (steel, wall conductance ratio $c$) are still missing; so far there are only
  perfectly insulating and perfectly conducting walls. For steel ducts the case in between is the most important one.
\item \textbf{Single GPU only}, and the field is uniform so far. The coil field ($\mathbf{B}(\mathbf{x})$ from
  Biot--Savart) would be the step towards coupling with the stellarator.
\end{itemize}
